\documentclass[aspectratio=169,11pt]{beamer}

\usepackage[brazilian]{babel}
\usepackage[utf8]{inputenc}
\usepackage[T1]{fontenc}
\usepackage{amsmath,amssymb,amsfonts,mathtools}
\usepackage{booktabs}
\usepackage{listings}
\usepackage{xcolor}

\usetheme{Madrid}
\setbeamertemplate{navigation symbols}{}

\lstset{
  basicstyle=\ttfamily\small,
  columns=fullflexible,
  keepspaces=true,
  showstringspaces=false,
  frame=single,
  breaklines=true
}

\title[Programação Funcional]{Programação Funcional:\\do Cálculo Lambda à Implementação de Métodos Numéricos}
\author[Ribas et al.]{Enzo Rocha Leite Diniz Ribas \and Luis A. D'Afonseca \and Lucas Martins Rocha}
\institute[CEFET-MG]{Centro Federal de Educação Tecnológica de Minas Gerais\\Belo Horizonte -- MG}
\date{Semana de Iniciação Científica -- CEFET-MG}

\begin{document}

\begin{frame}\titlepage\end{frame}

\begin{frame}{Roteiro}
    \tableofcontents
\end{frame}

\section{Motivação}

\begin{frame}{A pergunta central}
    \begin{center}\Large\textbf{O que significa implementar uma função matemática?}\end{center}
    \vspace{0.4cm}
    \begin{itemize}
        \item Na matemática, uma função relaciona entradas válidas a saídas.
        \item Na programação, uma função pode modificar estado, realizar E/S ou produzir
              erros.
        \item A programação funcional busca aproximar a computação do comportamento
              matemático das funções.
    \end{itemize}
    \begin{block}{Ideia-guia}
        Quanto mais explícito for o comportamento de uma função, mais fácil é raciocinar, testar e compor o programa.
    \end{block}
\end{frame}

\begin{frame}{Do problema matemático ao programa}
    Considere
    \[
        f(x)=\sqrt{x},\qquad \operatorname{Dom}(f)=\{x\in\mathbb{R}:x\geq0\}.
    \]
    \begin{itemize}
        \item $f(4)=2$
        \item $f(0)=0$
        \item $f(-1)$ não possui resultado em $\mathbb{R}$
    \end{itemize}
    \begin{alertblock}{Problema computacional}
        Como representar, no programa, que nem toda entrada possui um resultado válido?
    \end{alertblock}
\end{frame}

\section{Programação Imperativa}

\begin{frame}[fragile]{Uma implementação imperativa}
    \begin{lstlisting}[language=Java]
double sqrt(double x) {
    if (x < 0)
        throw new IllegalArgumentException();

    return Math.sqrt(x);
}
\end{lstlisting}
    \begin{itemize}
        \item comportamento normal;
        \item condição de erro;
        \item exceção durante a execução.
    \end{itemize}
\end{frame}

\begin{frame}{Exceções: uma solução, mas com custo conceitual}
    \begin{itemize}
        \item Exceções são importantes para representar falhas.
        \item O fluxo de erro fica separado do fluxo normal.
        \item A assinatura tradicional nem sempre comunica as falhas possíveis.
    \end{itemize}
    \begin{exampleblock}{Pergunta}
        E se o possível erro fizesse parte do próprio resultado?
    \end{exampleblock}
    \[
        \texttt{sqrt}:\mathbb{R}\longrightarrow\texttt{Maybe}\;\mathbb{R}
    \]
\end{frame}

\begin{frame}{Testes e contratos}
    \[
        \begin{array}{c|c}
            \text{Entrada} & \text{Comportamento esperado} \\ \hline
            4              & 2                             \\
            0              & 0                             \\
            -1             & \text{falha controlada}
        \end{array}
    \]
    \begin{itemize}
        \item Testes verificam comportamentos específicos.
        \item Tipos e contratos podem tornar erros explícitos.
        \item A programação funcional explora essa ideia sistematicamente.
    \end{itemize}
\end{frame}

\section{Programação Funcional}

\begin{frame}{O que é programação funcional?}
    A programação funcional trata \textbf{funções como abstrações fundamentais da computação}.
    \begin{itemize}
        \item funções de primeira classe;
        \item funções puras;
        \item imutabilidade;
        \item composição;
        \item recursão;
        \item efeitos colaterais controlados ou isolados.
    \end{itemize}
    \begin{block}{Função pura}
        Para uma mesma entrada, uma função pura produz sempre a mesma saída e não altera o estado observável do programa.
    \end{block}
\end{frame}

\begin{frame}{Função matemática $\times$ função pura}
    \[
        f(x)=x^2+1
    \]
    Uma função pura pode ser vista como
    \[
        f:\mathbb{R}\rightarrow\mathbb{R}.
    \]
    \begin{itemize}
        \item Não depende de variável global.
        \item Não altera variável externa.
        \item Não escreve em arquivo.
        \item Não depende de entrada externa.
    \end{itemize}
    \begin{block}{Resultado}
        A função pode ser analisada como uma transformação entrada $\rightarrow$ saída.
    \end{block}
\end{frame}

\begin{frame}[fragile]{Efeitos colaterais}
    \begin{lstlisting}
int counter = 0;

int f(int x) {
    counter++;
    return x * x;
}
\end{lstlisting}
    Embora pareça simples, o comportamento observável depende de estado externo.
    \begin{alertblock}{Consequência}
        A função deixa de ser apenas uma transformação entre entrada e saída.
    \end{alertblock}
\end{frame}

\section{Cálculo Lambda}

\begin{frame}{A base formal: Cálculo Lambda}
    O cálculo lambda fornece um modelo formal para representar funções e computações.
    \begin{enumerate}
        \item \textbf{Abstração}: definir uma função;
        \item \textbf{Aplicação}: aplicar uma função;
        \item \textbf{Substituição}: avaliar a aplicação.
    \end{enumerate}
    Exemplo:
    \[
        \lambda x.\;x^2
    \]
    representa uma função que recebe $x$ e retorna $x^2$.
\end{frame}

\begin{frame}{BNF: descrevendo a sintaxe}
    Uma gramática simples:
    \[
        \begin{aligned}
            \langle expr\rangle::={} & \langle var\rangle                                \\
                                     & \mid\lambda\langle var\rangle.\langle expr\rangle \\
                                     & \mid\langle expr\rangle\;\langle expr\rangle
        \end{aligned}
    \]
    \begin{itemize}
        \item $\lambda x.e$: abstração;
        \item $e_1e_2$: aplicação;
        \item $\langle var\rangle$: variável.
    \end{itemize}
    \begin{block}{Ponte conceitual}
        O cálculo lambda conecta a definição matemática de função a uma representação formal de computação.
    \end{block}
\end{frame}

\begin{frame}{Aplicação e redução}
    Considere
    \[
        (\lambda x.\;x*x)\;5.
    \]
    Pela substituição $x\leftarrow5$:
    \[
        5*5=25.
    \]
    A redução beta é
    \[
        (\lambda x.e)\;v\longrightarrow e[x:=v].
    \]
    \begin{block}{Ideia}
        Programar pode ser interpretado como construir e reduzir expressões.
    \end{block}
\end{frame}

\section{Currying e Composição}

\begin{frame}{Currying}
    Uma função pode ser escrita como
    \[
        f(x,y)=x+y.
    \]
    No cálculo lambda:
    \[
        f=\lambda x.\lambda y.\;x+y.
    \]
    Aplicação:
    \[
        f(2)(3)=5.
    \]
    \begin{block}{Currying}
        Uma função de vários argumentos é representada como uma sequência de funções de um único argumento.
    \end{block}
\end{frame}

\begin{frame}{Por que Currying é útil?}
    \[
        \operatorname{somar}\;x\;y=x+y.
    \]
    Fixando $x=10$:
    \[
        \operatorname{somar10}=\operatorname{somar}\;10.
    \]
    Logo,
    \[
        \operatorname{somar10}(5)=15.
    \]
    \begin{itemize}
        \item especialização;
        \item reutilização;
        \item composição;
        \item construção de pipelines.
    \end{itemize}
\end{frame}

\section{Monads e Efeitos}

\begin{frame}{Do erro para um valor}
    Em vez de
    \[
        \texttt{sqrt}(x)\rightarrow\text{exceção},
    \]
    podemos modelar
    \[
        \texttt{sqrt}(x)\rightarrow\texttt{Maybe}\;\mathbb{R}.
    \]
    \[
        \begin{cases}
            \texttt{Just}(y), & x\geq0, \\
            \texttt{Nothing}, & x<0.
        \end{cases}
    \]
    O possível fracasso passa a fazer parte do resultado.
\end{frame}

\begin{frame}{Monads: composição com contexto}
    Uma apresentação inicial de Monad:
    \begin{center}\textbf{uma abstração para compor computações que carregam um contexto.}\end{center}
    Exemplos:
    \begin{itemize}
        \item ausência: \texttt{Maybe};
        \item erro: \texttt{Either};
        \item entrada e saída: \texttt{IO};
        \item estado;
        \item listas.
    \end{itemize}
    \begin{block}{Cuidado conceitual}
        Monad não é simplesmente ``tratamento de exceções''. O ponto central é a composição de computações contextualizadas.
    \end{block}
\end{frame}

\begin{frame}{Compondo operações que podem falhar}
    Considere
    \[
        x\xrightarrow{\sqrt{\phantom{x}}}\sqrt{x}
        \xrightarrow{\log}\log(\sqrt{x}).
    \]
    Ambas possuem restrições de domínio. Com um contexto como \texttt{Maybe},
    podemos compor:
    \[
        \texttt{Maybe}\;A
        \longrightarrow
        (A\rightarrow\texttt{Maybe}\;B)
        \longrightarrow
        \texttt{Maybe}\;B.
    \]
    Se uma etapa falhar, a composição preserva o contexto de falha.
\end{frame}

\section{Métodos Numéricos}

\begin{frame}{Por que isso importa para métodos numéricos?}
    Métodos numéricos são algoritmos construídos a partir de transformações matemáticas.
    Para Newton--Raphson:
    \[
        x_{n+1}=x_n-\frac{f(x_n)}{f'(x_n)}.
    \]
    Definindo
    \[
        N_f(x)=x-\frac{f(x)}{f'(x)},
    \]
    temos
    \[
        x_{n+1}=N_f(x_n).
    \]
    \begin{block}{Perspectiva funcional}
        A iteração pode ser vista como aplicação repetida de uma função.
    \end{block}
\end{frame}

\begin{frame}{Newton-Raphson: formulação funcional}
    \[
        N(f,f')(x)=x-\frac{f(x)}{f'(x)}.
    \]
    A iteração:
    \[
        x_0,\quad N(x_0),\quad N^2(x_0),\quad N^3(x_0),\ldots
    \]
    ou
    \[
        x_{n+1}=N(x_n).
    \]
    \begin{itemize}
        \item $f$ é uma função recebida como argumento;
        \item $f'$ é uma função recebida como argumento;
        \item $N$ constrói a próxima aproximação.
    \end{itemize}
\end{frame}

\begin{frame}[fragile]{Newton-Raphson: versão imperativa}
    \begin{lstlisting}
x = x0;

for (int i = 0; i < maxIter; i++) {
    x = x - f(x) / df(x);

    if (abs(f(x)) < tol)
        break;
}
\end{lstlisting}
    O estado é atualizado:
    \[
        x^{(0)}\rightarrow x^{(1)}\rightarrow x^{(2)}
        \rightarrow\cdots
    \]
    \begin{alertblock}{Contraste}
        A implementação é descrita como uma sequência de mutações de estado.
    \end{alertblock}
\end{frame}

\begin{frame}[fragile]{Newton-Raphson: perspectiva funcional}
    \begin{lstlisting}[language=Haskell]
newtonStep f df x =
    x - f x / df x
\end{lstlisting}
    Sua estrutura é:
    \[
        (f,f',x)\longmapsto x-\frac{f(x)}{f'(x)}.
    \]
    Uma iteração completa pode ser construída pela aplicação repetida de
    \texttt{newtonStep}.
    \begin{block}{Observação}
        A estrutura matemática da iteração fica explícita no programa.
    \end{block}
\end{frame}

\begin{frame}{Casos problemáticos em métodos numéricos}
    Métodos numéricos possuem condições de falha:
    \begin{itemize}
        \item derivada nula;
        \item divergência;
        \item máximo de iterações;
        \item domínio inválido;
        \item perda de precisão numérica.
    \end{itemize}
    Uma representação explícita pode ser:
    \[
        \texttt{NewtonResult}=
        \texttt{Converged}\mid
        \texttt{MaxIterations}\mid
        \texttt{InvalidDerivative}\mid\cdots
    \]
    Assim, o algoritmo não precisa representar todos os resultados como um simples
    número.
\end{frame}

\section{Vantagens e Limitações}

\begin{frame}{Principais vantagens}
    \begin{columns}
        \column{0.48\textwidth}
        \textbf{Funções puras}
        \begin{itemize}
            \item previsibilidade;
            \item testes;
            \item raciocínio matemático.
        \end{itemize}
        \textbf{Imutabilidade}
        \begin{itemize}
            \item menos estado compartilhado;
            \item menos dependências ocultas.
        \end{itemize}
        \column{0.48\textwidth}
        \textbf{Composição}
        \begin{itemize}
            \item modularidade;
            \item pipelines.
        \end{itemize}
        \textbf{Tipos e contextos}
        \begin{itemize}
            \item erros explícitos;
            \item contratos mais claros.
        \end{itemize}
        \textbf{Currying}
        \begin{itemize}
            \item especialização;
            \item reutilização.
        \end{itemize}
    \end{columns}
\end{frame}

\begin{frame}{Programação funcional não elimina o imperativo}
    \begin{center}\Large\textbf{Paradigmas são ferramentas, não religiões.}\end{center}
    \vspace{0.3cm}
    Linguagens modernas combinam paradigmas:
    \begin{itemize}
        \item Java: lambdas e streams;
        \item Python: funções de primeira classe;
        \item JavaScript: composição e funções de alta ordem;
        \item Rust: imutabilidade, pattern matching e tipos algébricos;
        \item Haskell: forte orientação funcional.
    \end{itemize}
    \begin{block}{Questão de engenharia}
        A escolha deve considerar clareza, segurança, desempenho, domínio do problema e custo de desenvolvimento.
    \end{block}
\end{frame}

\section{Conclusão}

\begin{frame}{Conclusão}
    \begin{enumerate}
        \item Funções matemáticas fornecem uma base conceitual natural.
        \item O imperativo descreve computação por mudanças de estado.
        \item O funcional privilegia funções puras, imutabilidade e composição.
        \item O cálculo lambda fornece uma fundamentação formal para funções e aplicação.
        \item Currying permite especialização e composição.
        \item Monads estruturam computações que carregam contexto.
        \item Métodos numéricos podem ser expressos por composição e aplicação de funções.
    \end{enumerate}
\end{frame}

\begin{frame}{Mensagem final}
    \begin{center}
        \Huge\textbf{Do cálculo à computação.}
        \vspace{0.6cm}

        \Large
        Funções matemáticas
        $\rightarrow$
        Cálculo Lambda
        $\rightarrow$
        Programação Funcional
        $\rightarrow$
        Métodos Numéricos
    \end{center}
    \vfill
    \begin{center}\normalsize Obrigado.\end{center}
\end{frame}

\begin{frame}[allowframebreaks]{Referências}
    \footnotesize
    \begin{thebibliography}{9}

        \bibitem{pierce}
        PIERCE, B. C.
        \newblock \emph{Types and Programming Languages}.
        \newblock MIT Press, 2002.

        \bibitem{hutton}
        HUTTON, G.
        \newblock \emph{Programming in Haskell}.
        \newblock Cambridge University Press.

        \bibitem{burden}
        BURDEN, R. L.; FAIRES, J. D.
        \newblock \emph{Numerical Analysis}.
        \newblock Cengage Learning.

        \bibitem{abelson}
        ABELSON, H.; SUSSMAN, G.
        \newblock \emph{Structure and Interpretation of Computer Programs}.
        \newblock MIT Press.

    \end{thebibliography}
\end{frame}

\end{document}
